\section{\webappname}\label{sec:web-application} % * 6 pages, 7max

\paragraph{i will add smth about google analytics after i will implement it (We are using google analytics to know more about app usage)}

The main goal of \webappname{} is to let users without GIS expertise explore climate
data for any location in the world through accessible, easy-to-understand
visualisations. To this end, we designed it as an interactive map application, a common
solution for accessing geospatial data: users select a location on the map and obtain
its climate profile, which can then be compared with other locations or periods. The
design follows three principles: (i)~global coverage of climate data; (ii)~interactive
comparison across locations and periods; and (iii)~portability, achieved through web
technologies, with the browser acting as a universal client, and by keeping user
requirements to a minimum, without registration or local software installation. Climate data are retrieved
through an API rather than processed from files by the application itself
(Section~\ref{sec:web-application:data}), which keeps \webappname{} lightweight while
preserving the full resolution of the dataset.

% This section describes the design of \webappname{}. Section~\ref{sec:related:usage-example} introduces a real-world task 
% example from forest management, from which the requirements in Section~\ref{sec:web-application:requirements} are derived.
% Section~\ref{sec:web-application:data} explains how WorldClim data and other integrated sources, 
% Section~\ref{sec:web-application:architecture} presents the architecture of the web application, and 
% Section~\ref{sec:web-application:features} describes the main features, illustrated with the running example.

\subsection{Usage Example}\label{sec:web-application:usage-example}

To illustrate the capabilities of \webappname{}, we use a task from forest
management practice proposed by \climateexpert{}. Forest management plans
typically include an assessment of the local climate over the preceding 30-year
period, which guides decisions such as which tree species are suitable for a site.
As the climate shifts, a manager updating such a plan needs to know not only the
current climate of the site but also how it has changed and how it varies across
the managed area.

Consider a forest manager updating the plan for a stand near \examplecity{}.
The manager needs to:
(T1)~obtain the climate profile of the site for the most recent normal (1991--2020);
(T2)~quantify how the climate has changed relative to an earlier normal (1961--1990); and
(T3)~compare the site with an alternative location that may serve as a climate analogue for future conditions.

With raw WorldClim data, each of these tasks would require locating the relevant
raster files, downloading them and processing them with GIS tools.
Section~\ref{sec:web-application:features} shows how \webappname{} can perform all of these tasks
easily without requirements of guidelines or any GIS expertise. 

% * R1–Rn list of requirements for the app
\subsection{Requirements}\label{sec:web-application:requirements}
The requirements of \webappname{} arose from two complementary sources. First, we
inspected the existing tools reviewed in Section~\ref{sec:related:existing-tools},
identifying both must-have features already present in them, such as place search,
Walter--Lieth climographs and chart export, and their main limitations, such as the
lack of comparison across locations and periods, coarse spatial resolution and limited
search options. Second, the requirements were refined in a panel of stakeholders,
comprising software developers, data scientists and climate domain experts, which also
contributed the running example (Section~\ref{sec:web-application:usage-example}). The
resulting requirements are divided into functional and non-functional ones.

% \begin{description}
%   \setlength{\itemsep}{5pt}
%   \item \textbf{R1 --- Location search.} Select any location on land by name with any language; search by coordinates or a click on the map.
%   \item \textbf{R2 --- Data Visual Representation.} Display the climate data as a Walter--Lieth diagram, and a standard climograph on each page.
%   \item \textbf{R3 --- City climate profile.} Display the climate profile, with climograms, and summary statistics.
%   \item \textbf{R4 --- Location comparison.} Compare the climate of two locations on shared axes.
%   \item \textbf{R5 --- Period comparison.} Compare the climate of a location across 30-year normals and individual years.
%   \item \textbf{R6 --- Data export.} Export charts as images and data in raw formats.
%   \item \textbf{R7 --- Shareable URL.} Ability to share url with selected plots and filters
%   \item \textbf{R8 --- Portability.} Run in any modern web browser, without registration or local software installation.
%   \item \textbf{R9 --- Localization.} Provide a multilingual interface.
% \end{description}

\subsubsection{Functional requirements.}
\begin{description}
  \setlength{\itemsep}{3pt}
  \item[FR1 --- Location search.] Select any location on land by its name in any
    supported language, by coordinates or by clicking on the map (T1--T3).
  \item[FR2 --- Climate profile.] Display the climate of a location as a Walter--Lieth
    diagram or a standard climograph, with summary statistics (T1).
  \item[FR3 --- Period comparison.] Compare the climate of a location across 30-year
    normals and individual years (T2).
  \item[FR4 --- Location comparison.] Compare the climate of two locations on shared
    axes (T3).
  \item[FR5 --- Data export.] Export charts as images and data in raw formats.
  \item[FR6 --- Shareable views.] Encode the selected location, period and filters in
    the URL, so that any view can be shared with other users.
\end{description}

\subsubsection{Non-functional requirements.}
\begin{description}
  \setlength{\itemsep}{3pt}
  \item[NFR1 --- Portability.] Run in any modern web browser, without registration or
    local software installation.
  \item[NFR2 --- Response time.] Show search suggestions while the user types and
    display a climate profile within a few seconds, with repeated queries answered
    faster through caching. % TODO: replace with measured values
  \item[NFR3 --- Usability.] Be usable by people without GIS expertise
    (assessed in Section~\ref{sec:evaluation}).
  \item[NFR4 --- Localisation.] Provide a multilingual interface.
\end{description}

\subsection{Data Sources and Integration}\label{sec:web-application:data}
- 1 page?
\subsubsection{WorldClim Access through SCRAPI}\label{sec:web-application:data:scrapi-worldclim}
\webappname{} relies on WorldClim (Section~\ref{sec:related:worldclim}) as its data
source. As discussed there, the archive is distributed as 8{,}836 GeoTIFF files, about
85.6\,GB in total. An application handling these files directly would need to store the
whole archive, read the relevant raster for every request and locate the pixel that
corresponds to the selected place, which would tie its implementation to the layout
and format of the files.

As an alternative, we decided to consult the data through an API. \webappname{} uses
SCRAPI~\cite{scrapi}, a REST API that exposes the WorldClim archive as a set of
resources: grids, one for each spatial resolution; cells within a grid; rasters, one for
each variable and period; and pixels, the values of a raster in a given cell. Instead of
converting the rasters in advance, SCRAPI reads the GeoTIFF files at request time and
offers raster-oriented operations, such as finding the cell that contains a pair of
coordinates or extracting the monthly pixel values of a variable for a given period and
spatial resolution. SCRAPI also serves the four additional climate normals described in
Section~\ref{sec:related:worldclim}.

This decision has two main benefits. First, it decouples \webappname{} from the data
format: the application never handles files, so changes in the organisation or format
of the archive only affect the API, which protects the application against such
variations. Second, the API offers uniform access to WorldClim that is not tied to
\webappname{} and can therefore be reused by other applications interested in these
data. Moreover, since only the requested values are transferred, the application
remains lightweight.

% WorldClim provides gridded climate data covering the globe at a spatial resolution
% suitable for location-based climate analysis. The dataset contains monthly climate
% variables and long-term averages, allowing the application to present both climate
% normals and observations for individual years. Rather than downloading complete
% GeoTIFF archives to the browser, \webappname{} accesses the required values SCRAPI~\cite{scrapi}.

% SCRAPI exposes raster operations through a REST interface. For a point query, the
% application sends the selected coordinates together with the requested variable,
% period, and aggregation options. SCRAPI resolves the corresponding raster cells and
% returns the values in a machine-readable response. The same approach supports
% regional queries while keeping raster storage and processing on the server. This
% integration reduces the amount of data transferred to the client and enables the
% application to use WorldClim's spatial and temporal coverage without requiring GIS
% software or local data management.

\subsubsection{Search Flows}\label{sec:web-application:data:search-flows}
- mention about geonames, why and how we use them
- mention about usage of wikidata for coordinates search
- about map usage and logic with click on map (idk either its crutual to mention about map, but after all its also an integration)

\subsection{Architecture}\label{sec:web-application:architecture}
- 1.5 pages with architecture schema which takes half a page
- i will draw better schema for architecture and paste it here, below the schema i want to explain how regular flow of the app 
usage works (SCRAPI, Redis, Solr, Wikidata); Guess whould be good to tell about caching strategies, how solar works and why
we decided to implement solr insead of wikidata api; mention 
tests and pipelines; 
regarding ci/cd integration i dont know if there is any sense to write technical staff, I can just say
that we have mechanisms to check quality of the app before pushing it to server and ability automatically update city search 
functionality on prod after adding new locales

% * main features of the app, with screenshots
\subsection{Main Features}\label{sec:web-application:features}
- 4-5 pages (including screenshots); probably would be good in examples show exactly results from Prof Felipe tasks?

\subsubsection{City Climate}\label{sec:web-application:features:city-climate}
- two or three screenshots with city climate view (one top and two bottoms with 2 different climograms); 

\begin{figure}[t]
  \centering
  \figplaceholder{4cm}{City Climate page: map with location pin,
    raster cell outline and summary statistics}
  \caption{City Climate page for \examplecity{} (1991--2020 normal).}
  \label{fig:city-climate}
\end{figure}

\begin{figure}[t]
  \centering
  \begin{minipage}[t]{0.49\textwidth}
    \centering
    \figplaceholder{4cm}{Walter--Lieth diagram}
  \end{minipage}
  \begin{minipage}[t]{0.49\textwidth}
    \centering
    \figplaceholder{4cm}{Standard climograph}
  \end{minipage}\hfill
  \caption{Climate profile of \examplecity{} (1991--2020) as (a) a Walter--Lieth diagram and (b) a standard
    climograph}
  \label{fig:climographs}
\end{figure}

\subsubsection{Compare Cities}\label{sec:web-application:features:compare-cities}
- show comparison table and climogram, guess here could be wl diagrams side-by side, and

\begin{figure}[t]
  \centering
  \figplaceholder{5cm}{Compare Cities page: statistics table and wl diagrams side by side}
  \caption{Comparison of the climate normals between \examplecity{} and \examplecity{}.}
  \label{fig:compare-periods}
\end{figure}

\subsubsection{Compare Periods}\label{sec:web-application:features:compare-periods}
- here could be defalult climograms side by side or overlapped?

\begin{figure}[t]
  \centering
  \figplaceholder{5cm}{Compare Periods page: statistics table and overlapped chart}
  \caption{Comparison of the climate normals for \examplecity{}.}
  \label{fig:compare-periods}
\end{figure}

\subsubsection{Data Export}\label{sec:web-application:features:data-export}
- here i whould like to show exported picture from city climate page and raw csv export example

\begin{figure}[t]
  \centering
  \figplaceholder{5cm}{exported png of}
  \caption{exported png of city climate page for \examplecity{} }
  \label{fig:compare-periods}
\end{figure}

\subsubsection{Location search}\label{sec:web-application:features:location-search}
Users can locate a place in three ways: by typing its name, by entering a pair of
coordinates, or by clicking on the map (R1). Name search queries a Solr index of
over 12~million populated places from GeoNames~\cite{geonames}, with names available
in all supported languages. Suggestions appear while typing and are ranked by
population, and fuzzy matching based on the Damerau--Levenshtein edit
distance~\cite{damerau1964technique} tolerates typing errors, so a query such as
``Valldolid'' still returns Valladolid.

- also would be good to breafly mention about not localized sharable urls, and the number of localizations 